\documentclass[border=10pt]{standalone}
\usepackage{tikz}
\usepackage{amsmath}

\begin{document}
\begin{tikzpicture}[
    every node/.style={font=\sffamily},
    label/.style={font=\sffamily\small, text=black},
]

% Parameter t goes from 0 to 100
% Phase 1 (t=0-15):   Assembled camera
% Phase 2 (t=15-35):  Explode into parts
% Phase 3 (t=35-50):  All parts labeled, hold
% Phase 4 (t=50-60):  Camera body fades out
% Phase 5 (t=60-70):  Shutter fades out
% Phase 6 (t=70-80):  PCB fades out
% Phase 7 (t=80-100): Only lens + CMOS remain, show focal distance

% --- Compute animation parameters ---
% Explosion progress (0 to 1, during t=15..35)
\pgfmathsetmacro{\explode}{min(1, max(0, (\PARAM - 15) / 20))}

% Fade-out opacities
\pgfmathsetmacro{\bodyOpacity}{max(0, min(1, (60 - \PARAM) / 10))}
\pgfmathsetmacro{\shutterOpacity}{max(0, min(1, (70 - \PARAM) / 10))}
\pgfmathsetmacro{\pcbOpacity}{max(0, min(1, (80 - \PARAM) / 10))}

% Label appearance (fade in during t=30..38)
\pgfmathsetmacro{\labelOpacity}{min(1, max(0, (\PARAM - 30) / 8))}

% Final label appearance (focal distance, t=82..90)
\pgfmathsetmacro{\finalOpacity}{min(1, max(0, (\PARAM - 82) / 8))}

% Explosion offsets (vertical separation)
\pgfmathsetmacro{\bodyY}{2.8 * \explode}
\pgfmathsetmacro{\lensX}{-2.5 * \explode}
\pgfmathsetmacro{\shutterY}{1.0 * \explode}
\pgfmathsetmacro{\sensorY}{-1.0 * \explode}
\pgfmathsetmacro{\pcbY}{-2.8 * \explode}

% Slide toward center after parts removed (t=70..85)
\pgfmathsetmacro{\consolidate}{min(1, max(0, (\PARAM - 70) / 15))}
\pgfmathsetmacro{\lensSlide}{\lensX * (1 - 0.4*\consolidate)}
\pgfmathsetmacro{\sensorSlide}{\sensorY * (1 - \consolidate)}

% Colors
\definecolor{garnet}{HTML}{73000A}
\definecolor{atlantic}{HTML}{466A9F}
\definecolor{congaree}{HTML}{1F414D}
\definecolor{horseshoe}{HTML}{65780B}
\definecolor{warmgrey}{HTML}{676156}
\definecolor{darkgrey}{HTML}{363636}

% === CAMERA BODY (top-most when exploded) ===
\begin{scope}[shift={(0, \bodyY)}, opacity=\bodyOpacity]
    % Camera body shell
    \fill[darkgrey] (-1.8, -0.7) rectangle (2.2, 0.7);
    \fill[black!80] (-1.8, 0.4) rectangle (2.2, 0.7);
    % Viewfinder bump
    \fill[darkgrey] (0.8, 0.7) rectangle (1.6, 1.1);
    % Grip texture
    \foreach \y in {-0.5,-0.35,...,0.3} {
        \draw[black!40, line width=0.3pt] (-1.7, \y) -- (-1.2, \y);
    }
    % Shutter button
    \fill[black!60] (1.8, 0.7) circle (0.12);
    % Mode dial
    \fill[black!50] (-0.8, 0.7) rectangle (-0.3, 0.85);
    % Label (fades with body)
    \pgfmathsetmacro{\bodyLabelOp}{min(\labelOpacity, \bodyOpacity)}
    \node[label, opacity=\bodyLabelOp, anchor=west] at (2.5, 0) {\textbf{Camera Body}};
    \node[label, opacity=\bodyLabelOp, anchor=west, text=warmgrey, font=\sffamily\footnotesize] at (2.5, -0.35) {housing, controls,};
    \node[label, opacity=\bodyLabelOp, anchor=west, text=warmgrey, font=\sffamily\footnotesize] at (2.5, -0.6) {viewfinder};
\end{scope}

% === LENS ASSEMBLY (left when exploded) ===
\begin{scope}[shift={(\lensSlide, 0)}]
    % Lens barrel
    \fill[darkgrey!80] (-2.8, -0.6) rectangle (-1.8, 0.6);
    \fill[darkgrey!60] (-3.2, -0.5) rectangle (-2.8, 0.5);
    % Lens rings
    \draw[black!40, line width=0.8pt] (-2.0, -0.6) -- (-2.0, 0.6);
    \draw[black!40, line width=0.8pt] (-2.5, -0.55) -- (-2.5, 0.55);
    % Front element (glass)
    \fill[atlantic!30, opacity=0.6] (-3.2, -0.4) -- (-3.5, -0.35) -- (-3.5, 0.35) -- (-3.2, 0.4) -- cycle;
    \draw[atlantic!60, line width=0.5pt] (-3.5, -0.35) -- (-3.5, 0.35);
    % Lens elements inside (simplified)
    \draw[atlantic!40, line width=0.4pt] (-2.3, -0.45) arc (260:280:3);
    \draw[atlantic!40, line width=0.4pt] (-2.7, -0.4) arc (260:280:2.5);
    % Label
    \pgfmathsetmacro{\lensLabelOp}{min(\labelOpacity, 1)}
    \node[label, opacity=\lensLabelOp, anchor=east] at (-3.8, 0) {\textbf{Lens Assembly}};
    \node[label, opacity=\lensLabelOp, anchor=east, text=warmgrey, font=\sffamily\footnotesize] at (-3.8, -0.35) {glass elements,};
    \node[label, opacity=\lensLabelOp, anchor=east, text=warmgrey, font=\sffamily\footnotesize] at (-3.8, -0.6) {aperture, focus ring};
\end{scope}

% === SHUTTER MECHANISM ===
\begin{scope}[shift={(0, \shutterY)}, opacity=\shutterOpacity]
    % Shutter curtain frame
    \draw[warmgrey, line width=1pt] (-0.8, -0.45) rectangle (0.8, 0.45);
    % Shutter blades
    \fill[warmgrey!70] (-0.7, -0.02) rectangle (0.7, 0.02);
    \fill[warmgrey!50] (-0.6, -0.35) -- (-0.3, -0.02) -- (0.3, -0.02) -- (0.6, -0.35) -- cycle;
    \fill[warmgrey!50] (-0.6, 0.35) -- (-0.3, 0.02) -- (0.3, 0.02) -- (0.6, 0.35) -- cycle;
    % Opening
    \fill[white] (-0.15, -0.08) rectangle (0.15, 0.08);
    % Label (fades with shutter)
    \pgfmathsetmacro{\shutterLabelOp}{min(\labelOpacity, \shutterOpacity)}
    \node[label, opacity=\shutterLabelOp, anchor=west] at (1.2, 0) {\textbf{Shutter}};
    \node[label, opacity=\shutterLabelOp, anchor=west, text=warmgrey, font=\sffamily\footnotesize] at (1.2, -0.3) {exposure control};
\end{scope}

% === CMOS SENSOR ===
\begin{scope}[shift={(0, \sensorSlide)}]
    % Sensor substrate
    \fill[congaree!80] (-0.65, -0.45) rectangle (0.65, 0.45);
    % Pixel grid
    \foreach \x in {-0.55,-0.40,...,0.55} {
        \foreach \y in {-0.35,-0.20,...,0.35} {
            \fill[atlantic!50] (\x, \y) rectangle (\x+0.1, \y+0.1);
        }
    }
    % Sensor border
    \draw[congaree, line width=1.2pt] (-0.65, -0.45) rectangle (0.65, 0.45);
    % Wire bonds (small lines at edges)
    \foreach \x in {-0.5,-0.3,...,0.5} {
        \draw[garnet!60, line width=0.3pt] (\x, 0.45) -- (\x, 0.55);
        \draw[garnet!60, line width=0.3pt] (\x, -0.45) -- (\x, -0.55);
    }
    % Label
    \pgfmathsetmacro{\sensorLabelOp}{min(\labelOpacity, 1)}
    \node[label, opacity=\sensorLabelOp, anchor=west] at (1.2, 0) {\textbf{CMOS Sensor}};
    \node[label, opacity=\sensorLabelOp, anchor=west, text=warmgrey, font=\sffamily\footnotesize] at (1.2, -0.3) {image capture};
\end{scope}

% === PCB ===
\begin{scope}[shift={(0, \pcbY)}, opacity=\pcbOpacity]
    % PCB board
    \fill[horseshoe!60] (-1.5, -0.35) rectangle (1.8, 0.35);
    % Traces
    \draw[horseshoe!90, line width=0.4pt] (-1.2, 0.1) -- (-0.5, 0.1) -- (-0.5, -0.15) -- (0.3, -0.15);
    \draw[horseshoe!90, line width=0.4pt] (-1.0, -0.1) -- (-0.2, -0.1) -- (-0.2, 0.2) -- (0.8, 0.2);
    \draw[horseshoe!90, line width=0.4pt] (0.5, -0.2) -- (1.2, -0.2) -- (1.2, 0.05) -- (1.5, 0.05);
    % Components (small rectangles)
    \fill[darkgrey] (0.3, -0.2) rectangle (0.5, -0.05);
    \fill[darkgrey] (0.8, 0.05) rectangle (1.1, 0.2);
    \fill[darkgrey] (-0.8, -0.25) rectangle (-0.5, -0.1);
    % Chip
    \fill[black] (-0.1, -0.15) rectangle (0.25, 0.15);
    % Connector
    \fill[warmgrey] (1.5, -0.15) rectangle (1.8, 0.15);
    % Label (fades with PCB)
    \pgfmathsetmacro{\pcbLabelOp}{min(\labelOpacity, \pcbOpacity)}
    \node[label, opacity=\pcbLabelOp, anchor=west] at (2.1, 0) {\textbf{PCB / ISP}};
    \node[label, opacity=\pcbLabelOp, anchor=west, text=warmgrey, font=\sffamily\footnotesize] at (2.1, -0.3) {image processing,};
    \node[label, opacity=\pcbLabelOp, anchor=west, text=warmgrey, font=\sffamily\footnotesize] at (2.1, -0.55) {signal pipeline};
\end{scope}

% === FOCAL DISTANCE ANNOTATION (final phase) ===
\begin{scope}[opacity=\finalOpacity]
    % Draw focal distance arrow between lens and sensor
    \pgfmathsetmacro{\lensRight}{-1.8 + \lensSlide}
    \pgfmathsetmacro{\sensorLeft}{-0.65}
    \pgfmathsetmacro{\arrowY}{-0.85 + \sensorSlide}
    \draw[garnet, line width=1.5pt, <->] (\lensRight, \arrowY) -- (\sensorLeft, \arrowY);
    \pgfmathsetmacro{\midX}{(\lensRight + \sensorLeft) / 2}
    \node[font=\sffamily\bfseries, text=garnet] at (\midX, \arrowY - 0.35) {$f$};
    \node[font=\sffamily\footnotesize, text=warmgrey] at (\midX, \arrowY - 0.65) {focal length};

    % Title text
    \node[font=\sffamily\large\bfseries, text=garnet, anchor=south] at (0, 2.2) {Only these matter for depth estimation};
\end{scope}

% Fixed bounding box
\path (-6, -4.5) rectangle (6, 3.5);

\end{tikzpicture}
\end{document}
